\section{The chamber law and the reference data}\label{sec:vtm-data}

\subsection{Normalization, flux and Gaussian domination}

The chamber operator is the finite-range operator of
\cite[\S2]{ChamberLattice}; its bond argument is a phase.
It acts on four spin components and $48$ chamber labels per
site. The fermions are antiperiodic in each coordinate.
For a rectangular coordinate torus with $V$ sites and every
side at least four, put $r_q=2\pi/q$. The unnormalized integral
of a curvature observable $O$ is
\begin{equation}\label{eq:vtm-law}
 Z[O]=\sum_{\substack{z\in C^2(\Z)\\dz=0}}
  \int_{[0,r_q)^{3V}}
 O(dA+r_qz)\exp\bigl(-\tfrac12\|dA+r_qz\|^2\bigr)
 \det(D^{\rm ch}[qA]+\boldsymbol\mu)\,\dd A.
\end{equation}
Cochains and the integer sum are on this torus; norms and
pairings count positively oriented cells. The identification
$(A,z)\sim(A+r_qk,z-dk)$ makes the integral independent
of the chosen fundamental domain. The law is $Z[O]/Z[1]$.
The partition function is finite and positive for every real
$\boldsymbol\mu$, and the determinant gives nonnegative weights;
see \cite[\S2--3]{ChamberLattice}. Flux is the three integer
periods of $z$, ordered as $(23,13,12)$.
We retain every flux sector in \eqref{eq:vtm-law}.

The coordinate-torus Gaussian bound
\cite[Proposition~5.1]{ChamberLattice} is
\begin{equation}\label{eq:slab-gaussian-domination}
 \E^{\rm af}\exp\langle g,F^{\rm u}\rangle
       \le\exp(\|g\|^2/2)
\end{equation}
for every real torus two-cochain $g$.
It holds for both odd and even coordinate sides.
The law is invariant under lattice translations and
$F^{\rm u}\mapsto-F^{\rm u}$. In particular all odd
moments of a real linear observable vanish.

On $C(n)$, let $\gamma_n$ be the standard Gaussian measure
on the exact real two-cochains, and let $A(u)$ be the
mean-zero coclosed potential of $u$.
For $\theta\in[-\pi,\pi]^3$ put
$H_{\theta,\nu}=\theta_\nu/(qn)$. The matter factor is
\begin{equation}\label{eq:external-matter-factor}
 z_{r,n}=\int_{[-\pi,\pi]^3}\frac{\dd\theta}{(2\pi)^3}
     \int\det(D^{\rm ch}[q(A(u)+H_\theta)]
                     +\boldsymbol\mu)\,\dd\gamma_n(u)>0.
\end{equation}
The subscript $r$ denotes the regulator. The exact zero-flux
factorization of \cite[\S3.4]{ReferenceResponse} gives
\begin{align}
 Z_n^{(0)}[1]&=Z_n^{\rm ph}z_{r,n},
                     \label{eq:external-zero-factorization}\\
 Z_n^{\rm ph}&=r_q^{\,n^3+2}(2\pi)^{n^3-1}
                         \exp(-\mathcal L(n)/2).
                     \label{eq:external-photon-factor}
\end{align}
where $\mathcal L(n)$ is the scalar-Laplacian logarithmic
determinant in \cref{sec:vtm-defect}.
Thus the normalized zero-flux curvature law is the Gaussian
law weighted by the twist-averaged determinant in
\eqref{eq:external-matter-factor}. The factorization retains
the gauge and harmonic volume factors; they will contribute
to the partition defect.

\subsection{The source and cell contract}\label{sec:reference-cells}

The source identities and estimates in this subsection are
proved in \cite[\S2]{ReferenceResponse}. We state the needed
contract with its support and normalization.

Write $p_{\rm cell}=n_*/L$ and $w=p_{\rm cell}/2$.
The reference side $n_*$ and $p_{\rm cell}$ are odd,
so $w$ is half-integral. The number $2L$ of cells per
side is even. In unit-lattice coordinates the open cells
are $Q_c^\circ=w(c+(0,1)^3)$.

For a two-form $g$ let
\[
 g_{[\eta]}(\sigma)=\eta^{-3/2}
        \int_{Q_\eta(\sigma)}g_{\mu\nu}(y)\,\dd y,\qquad
 Q_\eta(\sigma)=\eta\bigl(x+(e_\mu+e_\nu)/2+[-1/2,1/2)^3\bigr),
\]
when $\sigma=(x;\mu\nu)$. The source of the introduction obeys
$d\widehat a^o=(f^\eta)_{[\eta]}$.
After periodizing, this gives
\[
 \langle d\widehat a,F^{\rm u}\rangle=F^\sharp(f^\eta).
\]
All $f^\eta$ are supported in $Q_\beta$ for $\eta\le1$, and
\begin{equation}\label{eq:cells-test-error}
 \|f^\eta-f\|_{H^s}\le\eta^2\|\beta\|_{H^{s+4}}
                       \quad(s\in\R).
\end{equation}
The variance converges with the quantitative estimate
\begin{equation}\label{eq:cells-variance-limit}
 \left|\sqrt{v_\eta}-\|f\|_2\right|
 \le21\eta|\beta|_{C^3}(D_\beta^{\rm buf})^{3/2}
                       +\eta^2\|\beta\|_{H^4}.
\end{equation}
Here $|\beta|_{C^3}$ is the maximum component derivative norm
through order three. These estimates follow from the centered
difference symbols and the cell-average approximation.

For the placed probe, the coordinate map $g_c$ carrying the
base cell to cell $c$ is, in coordinate $\mu$,
$x_\mu\mapsto x_\mu+2iw$ if $c_\mu=2i$, and
$x_\mu\mapsto(2i+2)w-x_\mu$ if $c_\mu=2i+1$.
Transport oriented cochains by this map before periodizing.
Denote the placed and reflected covering source by
$\widehat a_c^o$; its periodization is $\widehat a_c$.
On every carrying torus with $p_{\rm cell}\mid n$ and $n\ge n_*$,
\begin{equation}\label{eq:cells-orthogonality}
 \langle d\widehat a_c,d\widehat a_{c'}\rangle_n
                         =v_\eta\mathbf1_{\{c=c'\}}.
\end{equation}
The placement condition \eqref{eq:cells-fit} ensures that their
support boxes are strictly inside the open cells and that the
boxes enlarged by one embed in the torus.
Each polynomial $P(X_c)$ is consequently read by the same
local scalar integrand on every such torus.

Here is the scalar observable class used in the transfer,
with the conventions of \cite[\S4.1, (46)--(50)]{ChamberLattice}.
For a coordinate box $B$, let $B^{(j)}$ be its contained
oriented lattice $j$-cells. A scalar local integrand is a Borel
gauge-invariant function $G$ of the real bonds in $B^{(1)}$
such that, for some $C<\infty$ and $\epsilon_B>0$,
\[
 |G(\widehat A)|\le C\exp\left[
   (\tfrac12-\epsilon_B)
          \sum_{\sigma\in B^{(2)}}|d\widehat A(\sigma)|^2\right].
\]
Gauge invariance means $G(\widehat A-d\lambda)=G(\widehat A)$.
A local lift of $(A,z)$ is a real bond cochain $\widehat A$
on $B$ with $d\widehat A=F^{\rm u}$ and
$\exp(iq\widehat A(b))=\exp(iqA(b))$.
The cited construction proves existence and independence of
the lift for these evaluations.

For an open set $U$, let $\mathcal A(U)$ be the linear span
of products $\Phi(d\widehat A)\prod_jG_j(\widehat A|_{B_j})$,
where the boxes $B_j\subset U$ are pairwise disjoint and
$\Phi$ is a bounded Borel function of finitely many curvature
values in $U$. On a quotient with periods $\Gamma$ the boxes
must embed: each family $B_j+[-1,1]^3+\gamma$ is disjoint,
and all $B_j+\gamma$ for distinct $(j,\gamma)$ are disjoint.

For the present probe take a supporting coordinate box
$B_0\Subset Q_0^\circ$ as in \cite[\S2.2]{ReferenceResponse}
and $B_c=g_cB_0$. The local expression for $P(X_c)$ is
$P(\sum_{\sigma\subset B_c}d\widehat a_c^o(\sigma)
                                      d\widehat A(\sigma))$.
Cauchy--Schwarz bounds its argument by
$\sqrt{v_\eta}(\sum_{\sigma\subset B_c}|d\widehat A(\sigma)|^2)^{1/2}$,
so it has a bound
$C_P\exp(\frac14\sum_{\sigma\subset B_c}|d\widehat A(\sigma)|^2)$.
Thus it belongs to $\mathcal A(Q_c^\circ)$.
The reflection is antilinear on coefficients; it fixes the
coefficients of these real Wick polynomials.

Even cell indices give the precise translation formula
\begin{equation}\label{eq:cells-translation-observable}
 c(y)=\frac{2y}{p_{\rm cell}\eta}\in2\Z^3,\qquad
 X_{c(y)}=F^\sharp(\tau_{v+y}f^\eta),
 \quad y\in(R^*/L)\Z^3.
\end{equation}
Two such cells coincide on the torus exactly when the
translations agree modulo its physical periods. Lattice
translation invariance removes the common placement $v$
from their product expectation.

\subsection{The normalized source identity}

On the reference torus define $a(t)=e^{-1}\sum_ct_c\widehat a_c$ and
\begin{equation}\label{eq:external-reference-partition}
 \mathring\Phi(t)=z_{r,n_*}^{-1}
 \int\frac{\dd\theta}{(2\pi)^3}\int
 \det(D^{\rm ch}[q(A(u)+H_\theta+a(t))]
                         +\boldsymbol\mu_N)\,\dd\gamma_{n_*}(u).
\end{equation}
The $\theta$ integral is over $[-\pi,\pi]^3$.
The source-shift identity of
\cite[\S3.6, Proposition~3.2 and (53)]{ReferenceResponse},
valid for arbitrary real lattice mass, is
\begin{equation}\label{eq:external-Wick-product}
 \E^{(0)}_{n_*}\prod_c W_m(X_c;v_\eta,c_m)
       =\left[\prod_c(\partial_{t_c}^m-c_m)\right]\mathring\Phi(0),
       \qquad m=2,4.
\end{equation}
Thus \eqref{eq:vtm-reference-root} is exactly the response norm
controlled there. Its reference law is normalized, and its
coefficients are the ordinary free derivatives
\eqref{eq:response-coefficient-relations}. The vacuum parameter
of a common tuned trajectory cancels in this ratio. In a
two-volume logarithmic defect its extensive vacuum contributions
cancel as well.

